% Prepared and Customized by Dr. Mahdi Maktab Dar for the School of Computing and Information Science, ARU 

\documentclass[12pt, twoside]{report}
\renewcommand{\contentsname}{Table of Contents}
\usepackage[utf8]{inputenc}
\usepackage{graphicx}
\graphicspath{ {Images/} }
\usepackage{lipsum}
\usepackage[a4paper,width=150mm,top=25mm,bottom=30mm,bindingoffset=6mm]{geometry}
\usepackage{fancyhdr}
\usepackage{booktabs}
\usepackage{graphicx}

% Watermark
\usepackage{graphicx}
\usepackage{background}
\backgroundsetup{contents=\includegraphics{Images/wtrmrk.png}, scale=2, opacity=0.03, angle=90}

% https://en.wikibooks.org/wiki/LaTeX/Floats,_Figures_and_Captions
\usepackage{caption}
\usepackage{subcaption}

% References

\usepackage{natbib}
\bibliographystyle{agsm}


% Bold figure captions
\usepackage[labelfont=bf]{caption}


% Page number and header and footer management 

\pagestyle{fancy} 
\fancyhead{}
\fancyhf{} 
% Set the right side of the footer to be the page number
\fancyfoot[OR,ER]{\thepage}

\renewcommand{\headrulewidth}{0pt}
\renewcommand{\footrulewidth}{0pt} 

% https://en.wikibooks.org/wiki/LaTeX/Text_Formatting#Line_Spacing
\usepackage{setspace}
%\singlespacing
\onehalfspacing
%\doublespacing
%\setstretch{1.1}

%********************************************************************************* 
%********************************************************************************* 
%********************************************************************************* 

\begin{document}
\input{title page.tex}
\pagenumbering{roman}

%*********************************************************************************   
\chapter*{Declaration}
% This is just an example.
% Please customize the following paragraph accordingly. 

I, \textbf{Philip J. Fry}, declare that the work in this dissertation titled \textit{“Design and Development of a Wonderful System Using Extraordinary Methods to Improve Everything”} is carried out by me. This work has not been submitted to Anglia Ruskin University or any other educational institution for the award of a degree or educational qualification. I also declare that the information published in this dissertation has been obtained and presented in accordance with academic rules and ethical conduct. Any information obtained from other sources has been properly referenced. 

%*********************************************************************************   
%*********************************************************************************   

\chapter*{Acknowledgement}

% This is just an example, directly taken from some random online source. Replace and reword accordingly. 
This thesis work is dedicated to my husband, Dr. John Zoidberg, who has been a constant source of support and encouragement during the challenges of graduate school and life. I am truly thankful for having you in my life. This work is also dedicated to my parents, William and Karen Ryan, who have always loved me unconditionally and whose good examples have taught me to work hard for the things that I aspire to achieve.
%*********************************************************************************   
% Your ToC, list of figures, and Tables will be generated here automatically subject to proper labeling and referencing. 
\tableofcontents
\listoffigures
\listoftables
%*********************************************************************************  
\chapter*{Abstract}
% Use the following guidelines to shape your abstract: 
% A good abstract must be around 250 to 350 words in total
% Allocate ~4 lines for context background info about the topic
% Allocate ~3 lines for talking about the research gap which is the motivation of your study.
% Allocate ~5 lines for explaining the research method and contribution. 
% Allocate ~3 lines for explaining the major findings.
% Allocate ~3 lines for a conclusive statement (research impact and implication) 
% The abstract must be a single long paragraph (250 to 350 words). Citation is not required in the abstract. Make sure your abstract exhibits the aforementioned sub-components. Dedicate one full page to your abstract. 
% Remove the following "\lipsum[1-1]" (which generates dummy text) and replace it with your abstract.
\lipsum[1-1]
%*********************************************************************************  


\chapter{Introduction} 
\pagenumbering{arabic}
\thispagestyle{empty}

\section{Overview}
% Think of this section as a short story written for a layman and not a domain expert. Describes the context and background information around the topic to help readers understand the basics and current state of the literature/topic/research gap.
% Do use a pyramid structure to elaborate your story i.e. placing the most fundamental (general) information in the lead paragraphs of the story, and then arranging the remaining details, from most general to most specific, in the subsequent paragraphs. 
% Start with general arguments and narrow them down into details. This section must be heavily cited and referenced.
% Expected length for this section is around 1 to 2 pages. 
% Remove the following "\lipsum[2-4]" (which generates dummy text) and replace it with your abstract.
\lipsum[4-6]

%*********************************************************************************  
\section{Problem Background}

% Elaborate on the issues that are surrounding the state of affairs of your chosen topics. Identify what is wrong or can be done better with the regard to the chosen topic and how it impacts society. 

% Of course, some solutions exist for those issues that you have just mentioned. These are the literature studies.  Talk about the main gaps/shortcomings in the literature which derive your research motivations. Elaborate on the shortcomings and limitations of the existing solutions in a critical manner.  Categorize and group these shortcomings based on their context/source/driver/nature. 

% Then, propose a tentative rough solution (short paragraph – 3 lines ) to these shortcomings. this is your research hypothesis. Outline your research questions. These questions must be answered by the end of your research.

% The goal of this section is to make a room for your research. Essentially this section justifies the necessity of your research by outlining the shortcomings of the existing solution for issues that are surrounding the state of affairs of your chosen topics. 

% This section must be heavily cited and referenced with recent (last 3 to 5 years studies). 
% Expected length for this section is around 2 to 3 pages. 
% Remove the following "\lipsum[8-10]" (which generates dummy text) and replace it with your abstract.

\lipsum[3-6]

% This is how to insert a basic figure.
% Please refer to https://www.overleaf.com/learn/latex/Inserting_Images for more info.

\begin{figure}[t]
\includegraphics[width=1\textwidth]{Images/problem_background.png}
\centering
\caption{Process of elaborating problem background}
\label{fig:problem_background}
\end{figure}

% You can cite literature studies in various ways. 
% You can use \cite{} to cite literature within your sentence or alternatively you can use \citep{} to cite studies at the end of your paragraph. Please see the following example: 

Pre Deep Learning object detection techniques, such as Boosted Cascade by \cite{viola2001rapid}, Histograms of Oriented Gradients (HOG) by \cite{dalal2005histograms} and Deformable Part Models (DPM) by \cite{felzenszwalb2010cascade} were relatively inaccurate and unreliable for real-world applications \citep{papageorgiou2000trainable}.

%*********************************************************************************  
\section{Research Aim}
% Outlines your overall research goal. Think of the aim as an extended research title. It must explain what is going to be achieved and how it is going to be achieved.
% Keep your Research Aim high-level and avoid low-level technical details. It must be relatively understandable for a non-expert.  
% Expected length for this section is around 1 to 2 paragraphs. 
\lipsum[1-1]


%*********************************************************************************  
\section{Research Objectives}

% Objectives are self-contained, self-sufficient milestones that collectively contribute toward the fulfillment of your research goal. Describe your objectives through bullet point statements. you should define 2 to 4 objectives for your research. Checkout the following examples: 
\begin{itemize}
  \item To investigate possible adverse impacts of lossy compression algorithms on BLE audio transmission measured by SNR metric.
  \item To design and develop a working prototype of a lossless compression algorithm capable of working with available  BLE bandwidth. 
\end{itemize}

%*********************************************************************************  
\section{Research Scope}

% Consists of 2 implicit sub-sections including in-scope and out-of-the-scope.
% In scope outlines the inner boundary of your work. Parameters that you are considering for measurement and monitoring. Factors and aspects of the topic that you are caring about and considering covering in your research. 

% Out of the scope outlines parameters, aspects, and issues that are not the concern of your research and won’t be investigated. Think of it as a disclaimer clause.  
\lipsum[1-1]

%*********************************************************************************  
\section{Abridged Methodology}
% Summarize your CHAPTER 3 (Research Methodology) in 2 to 3  paragraphs. 
\lipsum[1-1]


%*********************************************************************************  
\section{Contribution}
% Outline the final products/findings of your research that contributes toward the fulfillment of the research gap (stated in section 1.2) at a holistic and high level. 
% Talk about your overall research impact and how it contributes to the state of affairs of your chosen topic expressed in section 1.2. 
\lipsum[1-1]

%*********************************************************************************  
%*********************************************************************************  
%*********************************************************************************  


\chapter{Literature Review} 
\thispagestyle{empty}
\section{Overview}
% Describes context and background information about the topic and restates the research gap to help readers understand the basics and current state of the literature/topic. Think of it as an icebreaker section for the main body of your literature review.  
\lipsum[1-1]
%********************************************************************************* 
\section{Handwriting Recognition}

% From this point onward, the heading and subheadings of this chapter depend on the nature of your research. 
% For example, if your research is about Handwriting Recognition your 2.2 sub-heading can be “Handwriting Recognition”. 

% Start this section with a taxonomical study. Investigate and elaborate on this section in a hierarchical order. Start with the main general topic and then dig into details including explaining different approaches, techniques, aspects/perspectives, etc. 

% Design a hierarchical taxonomy diagram (as below) and explain each branch and their sub-branches and their corresponding scientific works. 

% Provide pros/cons, and critical analysis for each branch and each study. Create a sub-sub heading for every branch and sub-branch. This section must be meticulously cited and referenced. 

% Remember, chapter 2 is NOT about you and your work and your idea. It is about the works of others so do not talk about your desired/employed method in this chapter at all. 



\begin{figure}[t]
\includegraphics[width=1\textwidth]{Images/taxonomy.png}
\centering
\caption{Handwriting recognition taxonomy} 
\label{fig:taxonomy}
\end{figure}

Figure \ref{fig:taxonomy} shows the handwriting recognition taxonomy.
\lipsum[1-1]

\section{Approach A}
\lipsum[1-1]

% This is how to generate tables in Latex. You can use https://www.tablesgenerator.com/ to generate latex code for your desired table.

\begin{table}[]
\centering
\caption{Impact of SMOTE sampling on model performance}
\begin{tabular}{@{}lccc@{}}
\toprule
\textbf{Class}             & \textbf{Precision} & \textbf{Recall} & \textbf{F-score} \\ \midrule
\textbf{1 (fraud)}         & 0.15               & 0.63            & 0.24             \\
\textbf{0 (legal)}         & 1.00               & 0.99            & 1.00             \\
\textbf{1 (fraud)}         & 0.84               & 0.58            & 0.68             \\
\textbf{0 (legal)}         & 0.99               & 0.99            & 1.00             \\ \bottomrule
\end{tabular}
\label{table:smote}
\end{table}



\subsection{Method 1}
Method 1 is about this and that…etc….etc. It uses this and that techniques…. etc..etc……blah..blah…blah.
\cite{dalal2005histograms} used this method to detect this and that….etc..etc……blah..blah…blah. Despite reasonable accuracy figures, this study has higher computational complexity compared to…etc…etc.

\subsection{Method 2}
Method 2 is about this and that. This method is primarily made to counter those issues. This method benefits  from this feature.  \cite{felzenszwalb2010cascade} used this method to detect this and that….etc… blah..blah…blah.

\section{Approach B}
\lipsum[1-1]
\subsection{Method 3}
Method 3 is about this and that. This method is primarily made to counter those issues. This method benefits  from this feature.  \cite{felzenszwalb2010cascade} used this method to detect this and that….etc… blah..blah…blah.

\section{Approach C}
\lipsum[1-1]
\subsection{Method 4}
Method 4 is about this and that. This method is primarily made to counter those issues. This method benefits  from this feature.  \cite{felzenszwalb2010cascade} used this method to detect this and that….etc… blah..blah…blah.
\subsection{Method 5}
Method 5 is about this and that. This method is primarily made to counter those issues. This method benefits  from this feature.  \cite{felzenszwalb2010cascade} used this method to detect this and that….etc… blah..blah…blah.

\section{Literature Review Summary}
% Provides a written summary/conclusion to your literature review. Also, summarize related work in a table like the following:
\lipsum[3-5]


% the following code is in fact literature review summary table. Use the following web apps to generate tables:
% https://www.tablesgenerator.com/
% https://www.latex-tables.com/

\begin{table}[]
\centering
\caption{Literature review summary }
\label{tab:my-table}
\resizebox{\textwidth}{!}{%
\begin{tabular}{@{}lllll@{}}
\toprule
\textbf{\#} & \textbf{Study}       & \textbf{Method Description}                                                                                        & \textbf{Pros}                                                            & \textbf{Cons}                                                                 \\ \midrule
\textbf{1}  & \textbf{John et al.} & \begin{tabular}[c]{@{}l@{}}This study uses method 1 to identify \\ malign cancer cells…blah blah blah\end{tabular} & \begin{tabular}[c]{@{}l@{}}High accuracy\\ Speed\end{tabular}            & \begin{tabular}[c]{@{}l@{}}Reliance on hyperspectral \\ cameras.\end{tabular} \\
\textbf{2}  & \textbf{Dave et al.} & \begin{tabular}[c]{@{}l@{}}This study uses method 2 to identify \\ malign cancer cells…blah blah blah\end{tabular} & High accuracy                                                            & \begin{tabular}[c]{@{}l@{}}High computational \\ complexity\end{tabular}      \\
\textbf{3}  & \textbf{Anna et al.} & \begin{tabular}[c]{@{}l@{}}This study uses method 1 to identify \\ malign cancer cells…blah blah bla\end{tabular}  & \begin{tabular}[c]{@{}l@{}}Compatible with \\ this and that\end{tabular} & Limited to those                                                              \\ \bottomrule
\end{tabular}%
}
\end{table}


%*********************************************************************************
%********************************************************************************* 
%********************************************************************************* 
\chapter{Methodology} 
\thispagestyle{empty}
\section{Overview}
Restate the research problem situation and your corresponding proposed solution concept in a holistic and high-level narrative. Think of this section as an icebreaker for chapter 3. 
%********************************************************************************* 

\section{Research Framework}
Create a technical pipeline (block diagram) that outlines the consecutive steps (processes) that needed to be taken to get the research aim and objectives delivered and the artifact developed. Explain each step (process) in detail. Describe the operations in consecutive chronological order. Provide enough detail that renders your work reproducible. Reproducibility is essential in scientific work. Readers should be able to redo your work just by reading this chapter. Figure \ref{fig:framework} shows an example of a research framework diagram. The following subheadings are just an example of how the research framework can be elaborated. Change them according to your framework and your research. 

\begin{figure}[h]
\includegraphics[width=1\textwidth]{Images/framework.png}
\centering
\caption{Overall research framework}  
\label{fig:framework}
\end{figure}

\subsection{Raw Data Acquisition} 
Populate accordingly. Explain all the necessary details, parameters, experimental setups, and components in detail.  
\subsection{Data Wrangling and Preprocessing} 
Populate accordingly. Explain all the necessary details, parameters, experimental setups, and components in detail.  
\subsection{Feature Analysis} 
Populate accordingly. Explain all the necessary details, parameters, experimental setups, and components in detail.  

%********************************************************************************* 
\section{Dataset}
Explain your data collection method (for studies with primary data collection) or Dataset Structure (for empirical studies that mainly rely on secondary data). Outline the statistical moments, sample types, histograms, missing values, data types, and distribution. training, evaluation, and testing sets.
%********************************************************************************* 
\section{Evaluation Metrics}
Explain your success/evaluation/measurement criteria. what factors will be measured and how.  Metrics like Accuracy, False positive, false negative, true positive, sensitivity, loss, confusion matrix, specificity, F-Score, etc.  Note: do not exhibit your results here. Just explain what metrics will be used to measure the proposed model performance. By the way, this is how to create equations in latex. you can create inline equations such as \(\sqrt{x^2+1}\) one. Also, you can address your equation parameters like $a$ and $x$. Equation numbering will be automatic. 
% for more info on equations please refer to https://www.overleaf.com/learn/latex/Mathematical_expressions 


\begin{equation}
\sqrt{x^2+1}
\end{equation}


%********************************************************************************* 
%********************************************************************************* 
%********************************************************************************* 
\chapter{Results} 
\thispagestyle{empty}
\section{Overview}
Categories your finding into sub-sections based on their natures and where they contribute and research questions they answer. Use tables and charts (bar charts, pie charts, histograms, scatter plots) generously. Interpret your charts and results. Explain trends, behaviors, and statistical moments. This chapter is all about figures and numbers and charts. Leave high-level explanations for the discussion chapter.  
\section{Heading Based on Research Deliverable}
% Your result with regards to deliverable 1 goes here. Customize the heading accordingly.


\section{Heading Based on Research Deliverable}
% Your result with regards to deliverable 2 goes here. Customize the heading accordingly. 


%********************************************************************************* 
%********************************************************************************* 
%********************************************************************************* 
\chapter{Discussion} 
\thispagestyle{empty} 
\section{Overview}
Categories your discussion into sub-chapters based on their natures and where they contribute and research questions they answer. Follow the same structure as your results chapter. Move away [Zoom-out] from the numbers and explain what those numbers mean in a more general and holistic scene and how they contribute toward the fulfillment of the research gap. Interpret your findings using high-level explanations. Compare your work with other studies. Outline your method's advantages and drawbacks in comparison to other studies \citep{sullivan1997video}. 
\section{Heading Based on Research Deliverable}
% Your discussion with regard to deliverable 1 goes here. Customize the heading accordingly. 


\section{Heading Based on Research Deliverable}
% Your discussion with regard to deliverable 2 goes here. Customize the heading accordingly. 


%********************************************************************************* 
%********************************************************************************* 
%********************************************************************************* 
\chapter{Conclusion}  
\thispagestyle{empty} 
Restate the gap, answer the research question based on your research findings and appraise your work. Criticize your work in the form of future work. 

%********************************************************************************* 
%********************************************************************************* 
%********************************************************************************* 
\bibliography{references.bib}
\addcontentsline{toc}{chapter}{Bibliography}
\thispagestyle{empty}

\end{document}